\section{Method}\label{sec:method}

\subsection{Training-free: ref-KV injection}\label{sec:method-tf}

At each denoising step we run one additional forward pass of the
reference image through a subset (``band'') of the transformer's layers
in \emph{capture} mode: at each active layer, take the post-QK-RMSNorm
key and the value for reference tokens and store them. In the real
generation's forward pass (\emph{inject} mode), the stored reference
key/value is concatenated onto the target sequence's key/value at the
same layers, before scaled-dot-product attention (Fig.~\ref{fig:mechanism}a).
Injection strength is a per-step taper schedule on the injected
\emph{value} only, not the key -- QK-RMSNorm makes any key scalar
multiple a no-op, so only value-scaling has an effect.

\paragraph{The position problem.} Reading the transformer's coordinate
construction directly shows every image slot's spatial axes are indexed
from the same origin; a reference token grid and a target token grid
therefore share a coordinate origin. Injecting the reference's own
post-RoPE key (\texttt{refpos=keep}) causes it to attend as if physically
located in the target's top-left sub-block, producing literal
copy-paste ghosting (Fig.~\ref{fig:mechanism}b). Re-centering the
reference's RoPE coordinates to the middle of the target grid removes
the ghosting but introduces a different failure: full-precision
repositioned content collides with the model's own freedom to choose a
new pose, producing anatomical distortion in roughly a third of samples.

\paragraph{Frequency-gated fix.} We re-apply RoPE to the reference key at
the target-centered coordinate, but null all but the lowest-$k$
frequency-pair channels on the spatial axes to the identity rotation
($k{=}8$, found by a small sweep). This conveys coarse spatial layout
(``this content sits left-of-center, upper-half'') without exact
pixel-position binding, avoiding the composition-collision failure of
full re-centering while retaining more geometric information than fully
position-free injection. Neither widening the injection band alone nor
applying frequency-gating alone at the original band improves over the
position-free baseline; only widening the band \emph{and} switching the
newly-added layers to frequency-gated injection wins on the development
set below -- we interpret early/middle layers as
more likely to encode coarse global structure, so feeding them
content-plus-coarse-position is a qualitatively different signal than
feeding late layers full content. A WD14 tagger extracts structural
attributes (hair/eye color, gender, accessories) into a generic,
non-instructional T2I caption, addressing attributes that are freely
``what'' rather than ``where'' -- orthogonal to the KV channel, which
cannot reliably address these regardless of position-gating.

\subsection{Training-light: band-LoRA}\label{sec:method-tl}

The training-free mechanism's structural failure mode -- degenerating
into copying reference pixels wholesale -- motivates a second,
independent mechanism: a standard LoRA (rank $r$, $\alpha{=}r$) on
\texttt{to\_q/to\_k/to\_v/to\_out.0} within a contiguous band of layers,
trained end-to-end to \emph{route} to the right identity information
rather than being handed it directly. A production recipe (band 10--25,
$r{=}8$, 3.9M / 0.06\% of backbone parameters) was tuned informally
before this work tested whether that specific depth range is actually
special, or just the first configuration that worked.
